\documentclass{fiche}
\metadonnees{8}{Ce qui a dû se passer}{Bloc 2 — Ce qui est possible : futur et modaux}

\begin{document}
\entete

\objectifs{%
\textbf{Langue} : \emph{must}, \emph{may}, \emph{might}, \emph{can't} ; le modal au passé
(\emph{+ have} + participe) ; lexique de l'enquête.\par
\textbf{Texte} : Arthur Conan Doyle, \og The Adventure of the Speckled Band \fg (1892).\par
\textbf{Durée} : 1\,h\,30 — exercices 35 min, texte et questions 30 min, version 25 min.}

%% ===================================================================
\exercice[13 min]{Le degré de certitude}

\rappel{\textbf{Rappel.} \emph{Must} : j'en suis presque sûr. \emph{Can't} : je suis presque
sûr du contraire. \emph{May} et \emph{might} : c'est possible. Avec \emph{have} + participe,
le jugement porte sur le passé.}

\consigne{Complétez avec \emph{must}, \emph{can't}, \emph{may} ou \emph{might}, suivis au
besoin de \emph{have} + participe passé.}

\begin{anglais}
\begin{enumerate}[label=\arabic*.,itemsep=2.2mm,leftmargin=*]
  \item There is mud on your sleeve: you \trou{must have come} in a dog-cart.
  \item She left home before six, so she \trou{must have started} in the dark.
  \item That \trou{can't be} her stepfather; he is a much taller man.
  \item The whistle \trou{may have been} the wind --- she is not sure.
  \item He \trou{can't have heard} us: he was asleep at the other end of the house.
  \item ``Whatever your reasons \trou{may be}, you are perfectly correct.''
  \item She looks exhausted; she \trou{must be} at the end of her strength.
  \item The metallic sound \trou{might have come} from the shutters.
  \item Nobody answered the bell: they \trou{must have gone} out.
  \item These marks \trou{can't be} old; they are perfectly fresh.
  \item The doctor \trou{may know} something, but he will not speak.
  \item You look wretched: you \trou{can't have slept} at all.
\end{enumerate}
\end{anglais}

\note{\emph{Must} : 1, 2, 7, 9. \emph{Can't} : 3, 5, 10, 12. \emph{May} / \emph{might} : 4,
6, 8, 11.}

\note[Deux règles]{La négation de \emph{must} au sens de la déduction n'est pas
\emph{must not} mais \emph{can't} : \og ce ne peut pas être \fg. Et \emph{must have} +
participe ne se traduit pas par \og il a dû \fg{} au sens de l'obligation mais par \og il a
forcément \fg, \og il a sûrement \fg{} : en français le même verbe sert aux deux, en anglais
non (\emph{he had to leave} = il a été obligé de partir ; \emph{he must have left} = il est
sûrement parti).}

%% ===================================================================
\exercice[10 min]{Reformuler}

\rappel{\textbf{Rappel.} Un modal remplace toute une proposition d'opinion : \emph{I am sure
that} $\rightarrow$ \emph{must} ; \emph{it is possible that} $\rightarrow$ \emph{may}.}

\consigne{Réécrivez chaque phrase avec un modal, sans changer le sens.}

\begin{anglais}
\begin{enumerate}[label=\arabic*.,itemsep=3mm,leftmargin=*]
  \item I am sure that he left early.
    \lignes{1}
    \corr{He must have left early.}
  \item It is possible that she heard a whistle.
    \lignes{1}
    \corr{She may have heard a whistle.}
  \item I am certain that this is not her handwriting.
    \lignes{1}
    \corr{This can't be her handwriting.}
  \item Perhaps the old doctor is lying.
    \lignes{1}
    \corr{The old doctor may be lying.}
  \item I am sure they have got rid of the evidence.
    \lignes{1}
    \corr{They must have got rid of the
    evidence.}
  \item It is impossible that a stranger came in that night.
    \lignes{1}
    \corr{A stranger can't have come
    in that night.}
  \item Perhaps he did not understand the question.
    \lignes{1}
    \corr{He may not have understood the
    question.}
  \item I am sure she is telling the truth.
    \lignes{1}
    \corr{She must be telling the truth.}
\end{enumerate}
\end{anglais}

\note{Le modal porte sur ce que \emph{je} pense, jamais sur ce qui s'est passé : c'est
pourquoi on ne peut pas le mettre au prétérit. Pour parler du passé, on garde le modal au
présent et on met l'infinitif au parfait (\emph{have} + participe). \emph{He must left} et
\emph{he musted leave} n'existent pas.}

%% ===================================================================
\exercice[12 min]{L'enquête}
\consigne{a) Associez chaque mot à sa traduction.}

\begin{multicols}{2}
\begin{enumerate}[label=\arabic*.,itemsep=1.2mm,leftmargin=*]
  \item a will \hfill \trou{f}
  \item an heir \hfill \trou{c}
  \item an estate \hfill \trou{j}
  \item a stepfather \hfill \trou{a}
  \item an inquiry \hfill \trou{h}
  \item a coroner \hfill \trou{l}
  \item to swear \hfill \trou{d}
  \item to deceive \hfill \trou{b}
  \item bewilderment \hfill \trou{i}
  \item a shutter \hfill \trou{e}
  \item a hinge \hfill \trou{k}
  \item to encompass \hfill \trou{g}
\end{enumerate}
\columnbreak
\begin{enumerate}[label=\alph*.,itemsep=1.2mm,leftmargin=*]
  \item un beau-père
  \item tromper, abuser
  \item un héritier
  \item jurer, attester sous serment
  \item un volet
  \item un testament
  \item cerner, entourer
  \item une enquête, une instruction
  \item la stupéfaction
  \item un domaine, une propriété
  \item un gond
  \item un coroner (magistrat enquêtant sur les morts suspectes)
\end{enumerate}
\end{multicols}

\note[Deux familles]{\emph{Step-} marque la parenté par remariage : \emph{stepfather},
\emph{stepmother}, \emph{stepsister}. À ne pas confondre avec \emph{-in-law}, qui marque
l'alliance par mariage (\emph{mother-in-law}, la belle-mère au sens de la mère du conjoint).
\emph{A will} est à la fois l'auxiliaire du futur, la volonté et le testament : les trois
sens tiennent dans l'idée de ce qu'on veut.}

\consigne{b) Complétez avec un mot de la liste.}
\begin{anglais}
\begin{enumerate}[label=\arabic*.,itemsep=2mm,leftmargin=*]
  \item At the \trou{inquiry} the \trou{coroner} asked her whether she could \trou{swear} to
    the whistle.
  \item Her \trou{stepfather} was the last \trou{heir} of a ruined family.
  \item She stared at Holmes in \trou{bewilderment}.
  \item Among the crash of the gale she may have been \trou{deceived}.
\end{enumerate}
\end{anglais}

%% ===================================================================
\newpage
\section{Texte}

\noindent\small\textit{Une jeune femme voilée se présente à Baker Street au petit matin ;
elle tremble, et ce n'est pas de froid.}\normalsize

\begin{texte}
``We shall soon set matters right, I have no doubt. You have come in by train this morning, I
see.''

``You know me, then?''

``No, but I observe the second half of a return ticket in the palm of your left glove. You
must have started early, and yet you had a good drive in a
\gl[a dog-cart]{dog-cart}{une carriole attelée}, along heavy roads, before you reached the
station.''

The lady gave a violent start and stared in \gl[bewilderment]{bewilderment}{la stupéfaction}
at my companion.

``There is no mystery, my dear madam,'' said he, smiling. ``The left arm of your jacket is
\gl[to spatter]{spattered}{éclabousser} with mud in no less than seven places. The marks are
perfectly fresh. There is no vehicle \gl[save]{save}{excepté, sinon} a dog-cart which throws
up mud in that way, and then only when you sit on the left-hand side of the driver.''

``Whatever your reasons may be, you are perfectly correct,'' said she. ``I started from home
before six, reached Leatherhead at twenty past, and came in by the first train to Waterloo.
Sir, I can stand this \gl[a strain]{strain}{une tension nerveuse} no longer; I shall go mad
if it continues. I have no one to turn to --- none, save only one, who cares for me, and he,
poor fellow, can be of little aid. I have heard of you, Mr. Holmes; I have heard of you from
Mrs. Farintosh, whom you helped in the hour of her \gl[sore]{sore}{cruel, pressant} need. It
was from her that I had your address. Oh, sir, do you not think that you could help me, too,
and at least throw a little light through the dense darkness which surrounds me? At present
it is out of my power to reward you for your services, but in a month or six weeks I shall be
married, with the control of my own income, and then at least you shall not find me
ungrateful.''

Holmes turned to his desk and, unlocking it, drew out a small case-book, which he consulted.

``Farintosh,'' said he. ``Ah yes, I recall the case; it was concerned with an opal
\gl[a tiara]{tiara}{un diadème}. I think it was before your time, Watson. I can only say,
madam, that I shall be happy to devote the same care to your case as I did to that of your
friend. As to reward, my profession is its own reward; but you are at liberty to
\gl[to defray]{defray}{couvrir, rembourser} whatever expenses I may be put to, at the time
which suits you best. And now I beg that you will lay before us everything that may help us
in forming an opinion upon the matter.''

``Alas!'' replied our visitor, ``the very horror of my situation lies in the fact that my
fears are so vague, and my suspicions depend so entirely upon small points, which might seem
\gl[trivial]{trivial}{insignifiant} to another, that even he to whom of all others I have a
right to look for help and advice looks upon all that I tell him about it as the fancies of a
nervous woman. He does not say so, but I can read it from his \gl[soothing]{soothing}{
apaisant} answers and \gl[averted]{averted}{détourné} eyes. But I have heard, Mr. Holmes,
that you can see deeply into the \gl[manifold]{manifold}{multiple} wickedness of the human
heart. You may advise me how to walk amid the dangers which
\gl[to encompass]{encompass}{cerner, entourer} me.''

``I am all attention, madam.''
\end{texte}
\source{Arthur Conan Doyle, \og The Adventure of the Speckled Band \fg, \emph{The Adventures
of Sherlock Holmes}, Londres, George Newnes, 1892.}

\partiel[15 min]{Questions}
\consigne{\en{Answer in English, in full sentences.}}

\begin{enumerate}[label=\arabic*.,itemsep=2mm,leftmargin=*]
  \item \en{What are the two things Holmes deduces, and from which clues?}
    \lignes{3}
    \corr{That she came by train, from the half of a return ticket in her left glove; and
    that she first drove in a dog-cart along heavy roads, from the fresh mud spattered in
    seven places on her left sleeve, which only a dog-cart throws up, and only on the
    driver's left.}
  \item \en{Why does Holmes explain his reasoning at once?}
    \lignes{2}
    \corr{Because she is frightened and has taken it for magic. Explaining calms her and
    shows her that he works from facts, which is what she needs before she can talk.}
  \item \en{What does she say about money, and how does Holmes answer?}
    \lignes{2}
    \corr{She cannot pay now but will be married within six weeks and will then have her own
    income. He answers that his profession is its own reward, and that she may cover his
    expenses whenever it suits her.}
  \item \en{Why has she not been believed until now?}
    \lignes{3}
    \corr{Her fears are vague and rest on small points that would seem trivial to anyone
    else. Even the man who ought to help her treats them as the fancies of a nervous woman.}
  \item \en{``He does not say so, but I can read it from his soothing answers and averted
    eyes.'' What kind of reader of signs is she?}
    \lignes{3}
    \corr{The same kind as Holmes: she reads what is not said, from details of behaviour.
    The difference is that nobody takes her readings seriously, which is precisely why she
    has come to a man whose trade it is.}
  \item \en{Find three details that show how urgent her situation is.}
    \lignes{2}
    \corr{She started before six in the dark; she says she can stand the strain no longer and
    will go mad; she has no one to turn to; she is trembling.}
\end{enumerate}

\note[Les modaux du texte]{\emph{You must have started early} : déduction sur le passé.
\emph{Whatever your reasons may be} : possibilité. \emph{Which might seem trivial to
another} : possibilité atténuée, plus prudente que \emph{may}. \emph{Do you not think that
you could help me} : \emph{could} de politesse. \emph{You may advise me} : ici \emph{may}
n'est pas la permission mais la capacité reconnue à l'autre.}

%% ===================================================================
\newpage
\section{Version}
\consigne{Traduisez le passage en français. Helen Stoner raconte la nuit où sa sœur Julia est
morte, deux ans plus tôt, quelques jours avant son mariage.}

\begin{version}
It was a wild night. The wind was howling outside, and the rain was beating and splashing
against the windows. Suddenly, amid all the \gl[hubbub]{hubbub}{le tumulte} of the
\gl[a gale]{gale}{un coup de vent, une tempête}, there burst forth the wild scream of a
terrified woman. I knew that it was my sister's voice. I sprang from my bed, wrapped a shawl
round me, and rushed into the corridor. As I opened my door I seemed to hear a low whistle,
such as my sister described, and a few moments later a
\gl[clanging]{clanging}{métallique, retentissant} sound, as if a mass of metal had fallen. As
I ran down the passage, my sister's door was unlocked, and revolved slowly upon its hinges. I
stared at it horror-stricken, not knowing what was about to issue from it. By the light of the
corridor-lamp I saw my sister appear at the opening, her face
\gl[blanched]{blanched}{blêmi} with terror, her hands \gl[to grope]{groping}{tâtonner} for
help, her whole figure swaying to and fro like that of a drunkard. I ran to her and threw my
arms round her, but at that moment her knees seemed to give way and she fell to the ground.
She writhed as one who is in terrible pain, and her limbs were dreadfully convulsed. At first
I thought that she had not recognised me, but as I bent over her she suddenly shrieked out in
a voice which I shall never forget, `Oh, my God! Helen! It was the band! The speckled band!'
There was something else which she would \gl[fain]{fain}{volontiers} have said, and she
stabbed with her finger into the air in the direction of the Doctor's room, but a fresh
convulsion seized her and choked her words.
\end{version}
\source{\emph{Ibid.}}

\lignes{22}

\corr{\textbf{Proposition de traduction.} La nuit était déchaînée. Le vent hurlait dehors, et
la pluie battait les vitres et s'y écrasait. Soudain, au milieu de tout le tumulte de la
tempête, éclata le cri sauvage d'une femme terrifiée. Je reconnus la voix de ma sœur. Je
bondis de mon lit, jetai un châle sur mes épaules et me précipitai dans le corridor. Comme
j'ouvrais ma porte, il me sembla entendre un sifflement grave, tel que ma sœur l'avait
décrit, et quelques instants plus tard un bruit métallique, comme si une masse de métal était
tombée. Tandis que je courais dans le couloir, la porte de ma sœur, déverrouillée, tourna
lentement sur ses gonds. Je la fixai, glacée d'horreur, sans savoir ce qui allait en sortir.
À la lueur de la lampe du corridor, je vis ma sœur paraître dans l'embrasure, le visage
blêmi de terreur, les mains tâtonnant à la recherche d'un secours, tout le corps oscillant
de-ci de-là comme celui d'un homme ivre. Je courus à elle et l'entourai de mes bras, mais à
cet instant ses genoux parurent se dérober et elle tomba à terre. Elle se tordait comme
quelqu'un qui souffre atrocement, et ses membres étaient affreusement convulsés. Je crus
d'abord qu'elle ne m'avait pas reconnue, mais comme je me penchais sur elle, elle poussa
soudain un cri d'une voix que je n'oublierai jamais : \og Oh, mon Dieu ! Helen ! C'était le
ruban ! Le ruban tacheté ! \fg{} Elle aurait voulu dire encore autre chose, et son doigt
poignarda l'air dans la direction de la chambre du docteur, mais une nouvelle convulsion la
saisit et étouffa ses paroles.}

\note[Choix de traduction 1]{\emph{there burst forth the wild scream} : inversion après
\emph{there} pour retarder le sujet et produire la surprise. Le français peut garder l'ordre
(\og éclata le cri \fg) : c'est un de ses rares alignements avec l'anglais.}

\note[Choix de traduction 2]{\emph{as if a mass of metal had fallen} : \emph{as if} +
pluperfect pour un fait supposé antérieur ; en français \og comme si \fg{} + plus-que-parfait,
jamais de conditionnel après \og si \fg.}

\note[Choix de traduction 3]{\emph{her face blanched with terror, her hands groping for
help, her whole figure swaying} : trois propositions absolues, sans verbe conjugué ni
préposition. Le français les rend par des groupes détachés en apposition, précédés de
l'article défini (\og le visage blêmi\ldots, les mains tâtonnant\ldots \fg), et non par
\og avec son visage \fg.}

\note[Choix de traduction 4]{\emph{the band} : \emph{a band} est à la fois le ruban, la bande
de tissu et la bande de gens. Toute la nouvelle joue sur cette ambiguïté, que le français ne
peut pas rendre : c'est une perte assumée, et le titre traditionnel (\og Le ruban
moucheté \fg) tranche dans un sens.}

\note[Choix de traduction 5]{\emph{she would fain have said} : \emph{fain} est un adverbe
archaïque (\og volontiers \fg) ; le tour entier équivaut à \og elle aurait voulu dire \fg. À
reconnaître, pas à imiter.}

%% ===================================================================
\partiel{Lexique à mémoriser}
\begin{itemize}[label=--,itemsep=0.8mm,leftmargin=*]\small
  \item \textbf{a will} — un testament ; \textit{et} la volonté\textit{, et l'auxiliaire du futur}
  \item \textbf{an heir} — un héritier ; \textit{le} h \textit{est muet, comme dans} hour, honest
  \item \textbf{an estate} — un domaine ; real estate \textit{« l'immobilier »}
  \item \textbf{a stepfather} — un beau-père (par remariage)
  \item \textbf{an inquiry} — une enquête ; \textit{verbe} to inquire
  \item \textbf{to swear} — jurer ; swore, sworn
  \item \textbf{to deceive} — tromper ; \textit{nom} deceit\textit{, et non} deception \textit{« déception » — faux ami}
  \item \textbf{a strain} — une tension ; to be under strain
  \item \textbf{to stand} — supporter ; I can't stand it
  \item \textbf{to turn to somebody} — se tourner vers quelqu'un, recourir à
  \item \textbf{a gale} — un coup de vent
  \item \textbf{to howl} — hurler ; \textit{formation imitative}
  \item \textbf{to spring} — bondir ; sprang, sprung
  \item \textbf{to rush} — se précipiter
  \item \textbf{a hinge} — un gond
  \item \textbf{to sway} — osciller, vaciller
  \item \textbf{to writhe} — se tordre
  \item \textbf{a limb} — un membre ; \textit{le} b \textit{est muet, comme dans} numb \textit{(fiche 4)}
  \item \textbf{to shriek} — pousser un cri perçant
  \item \textbf{speckled} — tacheté ; a speck \textit{« une tache, un grain »}
\end{itemize}

\end{document}
